\documentclass[11pt,a4paper]{article}

\usepackage[utf8]{inputenc}
\usepackage{amsmath}
\usepackage{amsfonts}
\usepackage{amssymb}
\usepackage{graphicx}
\usepackage{geometry}
\usepackage{listings}
\usepackage{xcolor}
\usepackage{hyperref}
\usepackage{multirow}

\definecolor{codegreen}{rgb}{0,0.6,0}
\definecolor{codegray}{rgb}{0.5,0.5,0.5}
\definecolor{codepurple}{rgb}{0.58,0,0.82}
\definecolor{backcolour}{rgb}{0.95,0.95,0.92}

\lstset{
    backgroundcolor=\color{backcolour},   
    commentstyle=\color{codegreen},
    keywordstyle=\color{blue},
    numberstyle=\tiny\color{codegray},
    stringstyle=\color{codepurple},
    basicstyle=\ttfamily\small,
    breakatwhitespace=false,         
    breaklines=true,                 
    captionpos=b,                    
    keepspaces=true,                 
    numbers=left,                    
    numbersep=5pt,                  
    showspaces=false,                
    showstringspaces=false,
    showtabs=false,                  
    tabsize=2,
    frame=single
}

\geometry{
    a4paper,
    total={170mm,257mm},
    left=25mm,
    right=25mm,
    top=25mm,
}

\begin{document}

\begin{center}
    \vspace*{1cm}
    {\large \textbf{ASSIGNMENT 2: ADVANCED HEURISTICS AND METAHEURISTICS FOR THE KNAPSACK PROBLEM}} \\
    \vspace{0.8cm}
    
    \textbf{Mustapha Muntassir} \\
    Filière : ASEDS \\
    Institut National des Postes et Télécommunications (INPT) \\
    \vspace{0.2cm}
    Date: 04-25-2026
\end{center}

\vspace{0.8cm}

\noindent \textbf{ABSTRACT}
\vspace{0.2cm}

\noindent This report explores the 0-1 Knapsack Problem using provided benchmark instances ($n=20$ and $n=100$). We present the mathematical formulation of the problem and establish an exact baseline using the IBM CPLEX solver. The performance of five greedy heuristics is evaluated to determine their scalability. Furthermore, this study introduces and implements advanced metaheuristic strategies, specifically Variable Neighborhood Descent (VND) and Variable Neighborhood Search (VNS), to explore the trade-offs between computational time and solution quality in high-dimensional search spaces. Results are consolidated in a comparative table to analyze the efficiency of each method as the problem size increases.

\vspace{0.8cm}

\section{Description of the Knapsack Problem}
The 0-1 Knapsack Problem is a classic combinatorial optimization challenge centered on constrained resource allocation. Given a set of $n$ items, each with a utility $u_i$ and a weight $w_i$, the goal is to maximize total utility without exceeding a capacity $W$. Each item is indivisible, represented by the binary constraint $x_i \in \{0, 1\}$. While small instances are trivial, the problem is NP-hard, requiring significant computational resources as the number of items increases.

\section{Mathematical Formulation}
The 0-1 Knapsack Problem is formulated as an Integer Linear Programming (ILP) model:

\subsection*{Parameters}
\begin{itemize}
    \item $n$: Total number of items available for selection.
    \item $u_i$: Utility associated with item $i$, where $u_i > 0$.
    \item $w_i$: Weight of item $i$, where $w_i > 0$.
    \item $W$: Total capacity of the knapsack.
\end{itemize}

\subsection*{Decision Variables}
\begin{equation*}
    x_i \in \{0, 1\} \quad \text{for } i = 1, \dots, n
\end{equation*}

\subsection*{Objective Function}
\begin{equation*}
    \max Z = \sum_{i=1}^{n} u_i x_i
\end{equation*}

\subsection*{Constraints}
\begin{equation*}
    \sum_{i=1}^{n} w_i x_i \le W
\end{equation*}

\section{Exact Solution using IBM CPLEX}
To evaluate the scalability of exact optimization, we utilize the IBM ILOG CPLEX Optimizer. This section documents the performance of the OPL formulation as the problem size increases.

\subsection{Exact Results for $n=20$ and $n=100$}
For the small-scale benchmark ($n=20, W=35$), the optimal value $Z=68$ was found in 0.03s. Stepping up to $n=100$ items with $W=1400$, the solver reached the global optimum instantaneously:
\begin{itemize}
    \item \textbf{Optimal Objective Value ($Z^*$):} 551
    \item \textbf{CPU Execution Time:} 0.08 seconds
\end{itemize}

\subsection{Performance Summary Matrix}
The following table consolidates the exact results obtained. The increase in CPU time confirms that while CPLEX is highly efficient for these scales, computational effort grows with $n$.

\begin{table}[htbp]
    \centering
    \begin{tabular}{|c|c|c|c|}
        \hline
        \textbf{Instance (n)} & \textbf{Capacity (W)} & \textbf{Optimal Value (Z)} & \textbf{CPU Time (s)} \\
        \hline
        20 & 35 & 68 & 0.03 \\
        \hline
        100 & 1400 & 551 & 0.08 \\
        \hline
    \end{tabular}
    \caption{Performance matrix for the CPLEX Exact Method on the provided instances.}
    \label{tab:cplex_matrix}
\end{table}
\section{Implementation and Scalability of Greedy Heuristics}

Since the Knapsack problem is NP-hard, approximate methods are essential. In this section, we evaluate five different greedy strategies across the provided instances to analyze their accuracy and computational efficiency.

\subsection{Pseudo-code for the 5 Greedy Heuristics}

\textbf{1. Max Utility First}
\begin{lstlisting}[language=Python, numbers=none, frame=none]
Sort items in descending order of utility u[i]
Initialize Current_Weight = 0, Total_Profit = 0
For each item i in sorted list:
    If Current_Weight + w[i] <= W:
        Current_Weight = Current_Weight + w[i]
        Total_Profit = Total_Profit + u[i]
Return Total_Profit
\end{lstlisting}

\textbf{2. Min Weight First}
\begin{lstlisting}[language=Python, numbers=none, frame=none]
Sort items in ascending order of weight w[i]
Initialize Current_Weight = 0, Total_Profit = 0
For each item i in sorted list:
    If Current_Weight + w[i] <= W:
        Current_Weight = Current_Weight + w[i]
        Total_Profit = Total_Profit + u[i]
Return Total_Profit
\end{lstlisting}

\textbf{3. Max Density First}
\begin{lstlisting}[language=Python, numbers=none, frame=none]
Calculate density ratio r[i] = u[i] / w[i] for all items
Sort items in descending order of density r[i]
Initialize Current_Weight = 0, Total_Profit = 0
For each item i in sorted list:
    If Current_Weight + w[i] <= W:
        Current_Weight = Current_Weight + w[i]
        Total_Profit = Total_Profit + u[i]
Return Total_Profit
\end{lstlisting}

\textbf{4. Hybrid Density}
\begin{lstlisting}[language=Python, numbers=none, frame=none]
Calculate hybrid ratio r[i] = (u[i]^2) / w[i] for all items
Sort items in descending order of hybrid ratio r[i]
Initialize Current_Weight = 0, Total_Profit = 0
For each item i in sorted list:
    If Current_Weight + w[i] <= W:
        Current_Weight = Current_Weight + w[i]
        Total_Profit = Total_Profit + u[i]
Return Total_Profit
\end{lstlisting}

\textbf{5. Reverse Density First}
\begin{lstlisting}[language=Python, numbers=none, frame=none]
Calculate density ratio r[i] = u[i] / w[i] for all items
Sort items in ascending order of density r[i]
Initialize Current_Weight = sum(w), Total_Profit = sum(u)
While Current_Weight > W:
    Current_Weight = Current_Weight - w[i]
    Total_Profit = Total_Profit - u[i]
    i = i + 1
Return Total_Profit
\end{lstlisting}

\subsection{Performance Analysis ($n=20$ and $n=100$)}
For the provided instances, greedy heuristics based on density ratios are highly effective. As shown in the experimental results, both \textbf{Max Density} and \textbf{Reverse Density} reached the optimal values of $68$ (for $n=20$) and $551$ (for $n=100$), effectively matching the CPLEX baseline with exceptionally low execution times.

\subsection{Comparative Results Matrix}
The following table summarizes the profit and execution time for each heuristic.

\begin{table}[htbp]
    \centering
    \small
    \begin{tabular}{|l|cc|cc|}
        \hline
        \multirow{2}{*}{\textbf{Heuristic}} & \multicolumn{2}{c|}{\textbf{n=20}} & \multicolumn{2}{c|}{\textbf{n=100}} \\
        & Profit & Time (s) & Profit & Time (s) \\
        \hline
        Max Utility First & 66 & 0.000048 & 551 & 0.000110 \\
        Min Weight First  & 54 & 0.000065 & 532 & 0.000055 \\
        Max Density       & 68 & 0.000051 & 551 & 0.000100 \\
        Hybrid ($u^2/w$)  & 66 & 0.000038 & 551 & 0.000092 \\
        Reverse Density   & 68 & 0.000059 & 551 & 0.000116 \\
        \hline
        \textbf{CPLEX (Optimal)} & \textbf{68} & \textbf{0.03} & \textbf{551} & \textbf{0.08} \\
        \hline
    \end{tabular}
    \caption{Comparison of Greedy Heuristics against CPLEX Exact Solutions.}
\end{table}

\subsection{Key Observations}
The greedy algorithms successfully approximate the solutions with near-zero execution times. The \textbf{Reverse Density First} and \textbf{Max Density} approaches proved to be the most robust approximate methods, offering optimal solutions for these test cases instantly
\section{Variable Neighborhood Descent (VND) Metaheuristic}

To improve upon the greedy initial solutions, we implemented a Variable Neighborhood Descent (VND) metaheuristic. VND explores multiple neighborhood structures sequentially to escape local optima that a single neighborhood search might not bypass.

\subsection{Neighborhood Structures}
We defined three distinct neighborhood structures to explore the solution space:
\begin{itemize}
    \item \textbf{$N_1$ (1-Flip):} Involves changing the status of a single item (inclusion or exclusion).
    \item \textbf{$N_2$ (2-Flip):} Involves flipping the status of two items simultaneously.
    \item \textbf{$N_3$ (Swap):} Exchanges an item currently in the knapsack with one that is not, maintaining or improving the utility while respecting capacity.
\end{itemize}

\subsection{VND Pseudocode}
\begin{lstlisting}[language=Python, numbers=none, frame=none]
Initialize x = [0, 0, ..., 0]
Define Neighborhoods: N = [N_1, N_2, N_3]
k = 0, k_max = 3

While k < k_max:
    Find the best neighbor y in N[k](x)
    If value(y) > value(x):
        x = y
        k = 0
    Else:
        k = k + 1

Return x
\end{lstlisting}

\subsection{VND Results and Analysis}
The VND algorithm was initialized with a null vector ($x = [0]*n$) and executed until a local optimum was reached across all defined neighborhoods.

\begin{table}[htbp]
    \centering
    \begin{tabular}{|l|c|c|c|}
        \hline
        \textbf{Instance} & \textbf{Value (Z)} & \textbf{Weight (W)} & \textbf{Execution Time (ms)} \\
        \hline
        n=20  & 66  & 35   & 2.2997 \\
        \hline
        n=100 & 551 & 1398 & 310.9344 \\
        \hline
    \end{tabular}
    \caption{Performance results for the VND Metaheuristic.}
    \label{tab:vnd_results}
\end{table}

\subsection{Observations}
The results demonstrate that VND is highly effective at reaching the global optimum for the $n=100$ instance. While the execution time is significantly higher than the greedy heuristics (approximately 310 ms vs 0.1 ms), the trade-off is justified by the increase in solution quality. The algorithm successfully converged to the same optimal value provided by the IBM CPLEX solver.

\section{Variable Neighborhood Search (VNS) Metaheuristic}

Building upon the VND algorithm, we implemented a Variable Neighborhood Search (VNS) metaheuristic. While VND is effective, it can still become trapped in local optima depending on the initial solution. VNS introduces a \textit{shaking} mechanism to systematically escape these local optima before applying the VND local search.

\subsection{Shaking Function}
The shaking function acts as a perturbation phase. Given an incumbent solution $x$ and a neighborhood index $k$, it randomly flips $k$ items in the knapsack. To maintain feasibility, if a flip causes the total weight to exceed the capacity $W$, the flip is reverted. This ensures the search explores different areas of the solution space without spending computational effort on infeasible solutions.

\subsection{VNS Pseudocode}
\begin{lstlisting}[language=Python, numbers=none, frame=none]
Initialize best = [0, 0, ..., 0]
Define Local Search Neighborhoods: N = [N_1, N_2, N_3]
k_max = 3

While not termination_condition:
    k = 1
    While k <= k_max:
        # Shaking phase
        x_prime = shaking(best, k)
        
        # Local search phase
        x_second = N[k](x_prime)
        
        # Move or not
        If value(x_second) > value(best) and weight(x_second) <= W:
            best = x_second
            k = 1
        Else:
            k = k + 1

Return best
\end{lstlisting}

\subsection{VNS Results and Analysis}
The VNS algorithm was run with a predefined time limit of 100 seconds to evaluate its capacity to find optimal solutions within a bounded computational budget.

\begin{table}[htbp]
    \centering
    \begin{tabular}{|l|c|c|c|}
        \hline
        \textbf{Instance} & \textbf{Value (Z)} & \textbf{Weight (W)} & \textbf{Execution Time (s)} \\
        \hline
        n=20  & 67  & 35   & 100.0065 \\
        \hline
        n=100 & 551 & 1398 & 100.1540 \\
        \hline
    \end{tabular}
    \caption{Performance results for the VNS Metaheuristic.}
    \label{tab:vns_results}
\end{table}

\subsection{Observations}
The VNS approach found a near-optimal solution of 67 for the $n=20$ instance and successfully matched the global optimum of 551 for the $n=100$ instance. The execution time was capped by the termination condition (100 seconds), providing ample time for the algorithm to explore various neighborhood structures through the systematic shaking process.

\section{Comprehensive Comparative Analysis}
In this section, we compare the performance of all implemented methods. The results are categorized into Exact (CPLEX) and Approximate (Greedy, VND, and VNS) methods to analyze the trade-offs between optimality and computational speed.

\begin{table}[htbp]
    \centering
    \footnotesize
    \begin{tabular}{|l|cc|cc|cc|cc|cc|}
        \hline
        \multirow{3}{*}{\textbf{Instance}} & \multicolumn{2}{c|}{\textbf{Exact}} & \multicolumn{8}{c|}{\textbf{Approximate Methods}} \\
        \cline{2-11}
        & \multicolumn{2}{c|}{\textbf{CPLEX}} & \multicolumn{2}{c|}{\textbf{Greedy (Best)}} & \multicolumn{2}{c|}{\textbf{Other Greedy}} & \multicolumn{2}{c|}{\textbf{VND}} & \multicolumn{2}{c|}{\textbf{VNS}} \\
        \cline{2-11}
        & $Z$ & CPU (s) & $Z$ & CPU (s) & $Z^*$ & CPU (s) & $Z$ & CPU (s) & $Z$ & CPU (s) \\
        \hline
        \textbf{n=20} & 68 & 0.0300 & 68 & <0.0001 & 54-66 & <0.0001 & 66 & 0.0023 & 67 & 100.006 \\
        \hline
        \textbf{n=100} & 551 & 0.0800 & 551 & <0.0002 & 532 & <0.0002 & 551 & 0.3109 & 551 & 100.154 \\
        \hline
    \end{tabular}
    \caption{Summary Table: Exact vs Approximate Methods.}
\end{table}

\subsection{Detailed Breakdown of Greedy Heuristics}
To fulfill the requirement of comparing all specific greedy strategies, the following matrix details the performance of each individual heuristic:

\begin{table}[htbp]
    \centering
    \small
    \begin{tabular}{|l|cc|cc|}
        \hline
        \multirow{2}{*}{\textbf{Greedy Heuristic}} & \multicolumn{2}{c|}{\textbf{n=20 (W=35)}} & \multicolumn{2}{c|}{\textbf{n=100 (W=1400)}} \\
        \cline{2-5}
        & Value ($Z$) & Time (s) & Value ($Z$) & Time (s) \\
        \hline
        Max Utility First & 66 & 0.00004 & 551 & 0.00011 \\
        Min Weight First  & 54 & 0.00006 & 532 & 0.00005 \\
        Max Density       & \textbf{68} & 0.00005 & \textbf{551} & 0.00010 \\
        Hybrid ($u^2/w$)  & 66 & 0.00003 & 551 & 0.00009 \\
        Reverse Density   & \textbf{68} & 0.00005 & \textbf{551} & 0.00011 \\
        \hline
    \end{tabular}
    \caption{Detailed performance comparison of all 5 Greedy Heuristics.}
    \label{tab:greedy_results}
\end{table}

The comparative analysis reveals that for these specific instances, the greedy algorithms based on Density achieve the optimal value identified by CPLEX near-instantaneously. The VND approach provides robust results, converging to the optimal value for the larger instance in less than half a second. The VNS metaheuristic, bounded by a 100-second execution limit, also successfully locates the optimal/near-optimal bounds, demonstrating its capability for extensive solution space exploration. While CPLEX remains the benchmark for these scales, the approximate and metaheuristic approaches prove their validity and efficiency.

\section{Generalization: The Multidimensional Knapsack Problem (MKP)}
The 0-1 Knapsack Problem can be naturally extended to the Multidimensional Knapsack Problem (MKP), where items consume multiple resource constraints (e.g., volume, weight, budget) simultaneously rather than just a single capacity constraint.

\subsection{Mathematical Formulation of MKP}
The MKP is formulated by extending the capacity constraint into a system of $m$ constraints:
\begin{align*}
\text{Maximize} \quad & Z = \sum_{i=1}^{n} u_i x_i \\
\text{Subject to} \quad & \sum_{i=1}^{n} w_{ij} x_i \leq W_j \quad \forall j \in \{1, 2, \dots, m\} \\
& x_i \in \{0, 1\} \quad \forall i \in \{1, 2, \dots, n\}
\end{align*}
where $w_{ij}$ is the requirement of item $i$ for resource $j$, and $W_j$ is the total available capacity of resource $j$.

\subsection{Generalizing the Solution Approaches}
\textbf{IBM CPLEX:} Generalizing the exact method to the MKP is trivial when using a solver like CPLEX. The formulation simply requires adding multiple linear constraints rather than a single one. CPLEX handles the multi-constraint branch-and-bound optimization natively.

\textbf{Greedy Heuristics:} For greedy approaches, the concept of "density" ($u_i / w_i$) becomes ambiguous when there are multiple weights. To generalize, one must construct a pseudo-utility or surrogate weight that aggregates the multiple constraints. A common approach is dividing the utility by a weighted sum of the resources:
$ \text{Surrogate Ratio}_i = \frac{u_i}{\sum_{j=1}^{m} \lambda_j w_{ij}} $
where $\lambda_j$ are multipliers reflecting the scarcity or importance of constraint $j$.

\textbf{Local Search, VND, and VNS:} The core mechanisms of the metaheuristics (1-Flip, 2-Flip, Swap, Shaking) remain identical. The only modification required is in the \textit{feasibility check}. Instead of validating $\text{weight}(y) \leq W$, the algorithm must validate that the modified solution satisfies \textit{all} $m$ constraints: $ \text{weight}_j(y) \leq W_j \quad \forall j$. As long as the perturbation or local step results in a feasible vector across all dimensions, the metaheuristic frameworks function seamlessly.

\begin{thebibliography}{9}

\bibitem{MountassirRepo}
M. Mountassir.
\textit{0-1 Knapsack Problem Solver: Exact and Heuristic Methods}.
GitHub Repository, 2026. 
\url{https://github.com/Mount1ssir/-0-1_Knapsack_Problem}
\end{thebibliography}

\newpage
\appendix
\section{Python Source Code for Greedy Heuristics}
\begin{lstlisting}[language=Python]
# first algo 
def max_Value_First(W, w, u):
    items = sorted(zip(u, w), key=lambda x: x[0], reverse=True)
    s_u, s_w = zip(*items)
    c = 0
    p = 0
    l=[]
    for i in range(len(s_u)):
        if c + s_w[i] <= W:
            c += s_w[i]  
            p += s_u[i]
            l=l+[s_u[i]]
    return f"solution : {p}\nchosen items : {l}\nthe weight : {c}"
# second algo
def min_Value_First(W, w, u):
    items = sorted(zip(u, w), key=lambda x: x[1])
    s_u, s_w = zip(*items)
    c = 0
    p = 0
    l=[]
    for i in range(len(s_u)):
        if c + s_w[i] <= W:
            c += s_w[i]  
            p += s_u[i]
            l=l+[s_u[i]]
    return f"solution : {p}\nchosen items : {l}\nthe weight : {c}"
# third algo
def max_Density(W, w, u):
    r = []
    for i in range(len(w)):
        r = r + [u[i]/w[i]]  
    items = sorted(zip(r, u, w), key=lambda x: x[0], reverse=True)
    R, s_u, s_w = zip(*items)
    c = 0 # Current Weight
    p = 0 # Total Profit
    l = [] # List of chosen utilities
    for i in range(len(R)):
        if c + s_w[i] <= W:
            c += s_w[i]
            p += s_u[i]
            l = l + [s_u[i]]          
    return f"solution : {p}\nchosen items : {l}\nthe weight : {c}"
# fourth algo
def H_Density(W, w, u):
    r = []
    for i in range(len(w)):
        r = r + [(u[i]**2)/w[i]]
    items = sorted(zip(r, u, w), key=lambda x: x[0], reverse=True)
    R, s_u, s_w = zip(*items)
    c = 0 # Current Weight
    p = 0 # Total Profit
    l = [] # List of chosen utilities
    for i in range(len(R)):
        if c + s_w[i] <= W:
            c += s_w[i]
            p += s_u[i]
            l = l + [s_u[i]]
    return f"solution : {p}\nchosen items : {l}\nthe weight : {c}"
# fifth algo
def reverse_Density_First(W, w, u):
    r = []
    for i in range(len(w)):
        r = r + [u[i]/w[i]]
    
    items = sorted(zip(r, u, w), key=lambda x: x[0])
    R, s_u, s_w = zip(*items)
    
    current_weight = sum(w)
    total_profit = sum(u)
    in_knapsack = [True] * len(s_u)
    i = 0
    while current_weight > W and i < len(s_w):
        current_weight -= s_w[i]
        total_profit -= s_u[i]
        in_knapsack[i] = False
        i += 1

    s=[]
    for i in range (len(s_u)) : 
        if in_knapsack[i] : 
            s=s+[s_u[i]]
            
    return f"solution : {total_profit}\nchosen items : {s}\nthe weight : {current_weight}"
\end{lstlisting}

\section{Python Source Code for VND}
\begin{lstlisting}[language=Python]
def get_weight(x, n, w):
    return sum(w[i]*x[i] for i in range(n))

def get_value(x, n, u):
    return sum(u[i]*x[i] for i in range(n))

# 1-Flip
def local_search_1flip(x, n, W, u, w):
    while True:
        best_flip = -1
        max_gain = 0
        for i in range(n):
            y = list(x)
            y[i] = 1 - x[i]
            if get_weight(y, n, w) <= W:
                gain = get_value(y, n, u) - get_value(x, n, u)
                if gain > max_gain:
                    max_gain = gain
                    best_flip = i
        if best_flip != -1:
            x[best_flip] = 1 - x[best_flip]
        else:
            break
    return x

# 2-Flip
def local_search_2flip(x, n, W, u, w):
    while True:
        best_pair = (-1, -1)
        max_gain = 0
        for i in range(n):
            for j in range(i+1, n):
                y = list(x)
                y[i] = 1 - x[i]
                y[j] = 1 - x[j]
                if get_weight(y, n, w) <= W:
                    gain = get_value(y, n, u) - get_value(x, n, u)
                    if gain > max_gain:
                        max_gain = gain
                        best_pair = (i, j)
        if best_pair != (-1, -1):
            i, j = best_pair
            x[i] = 1 - x[i]
            x[j] = 1 - x[j]
        else:
            break
    return x

# Neighborhood 3 : Swap
def local_search_swap(x, n, W, u, w):
    x = list(x)
    while True:
        improvement = False
        x_best = list(x)
        for i in range(n - 1):
            for j in range(i + 1, n):
                if x[i] != x[j]:
                    y = list(x)
                    y[i], y[j] = y[j], y[i]
                    if get_weight(y, n, w) <= W and get_value(y, n, u) > get_value(x_best, n, u):
                        x_best = y
                        improvement = True
        if improvement:
            x = x_best
        else:
            break
    return x

# VND : Algorithm 9
def VND(n, W, u, w):
    x = [0]*n
    neighborhoods = [local_search_1flip, local_search_2flip, local_search_swap]
    k_max = len(neighborhoods)
    k = 0
    while k < k_max:
        y = neighborhoods[k](list(x), n, W, u, w)
        if get_value(y, n, u) > get_value(x, n, u):
            x = y
            k = 0
        else:
            k += 1
    return x
\end{lstlisting}

\section{Python Source Code for VNS}
\begin{lstlisting}[language=Python]
import random
import time

def shaking(x, k, n, W, w):
    y = list(x)
    for _ in range(k):
        i = random.randint(0, n - 1)
        y[i] = 1 - y[i]
        if get_weight(y, w) > W:
            y[i] = 1 - y[i] 
    return y

def VNS(n, W, u, w, time_limit=5):
    neighborhoods = [local_search_1flip, local_search_2flip, local_search_swap]
    k_max = len(neighborhoods)
    
    best = [0] * n 
    start = time.perf_counter()

    while (time.perf_counter() - start) < time_limit:
        k = 0
        while k < k_max:
            x_prime = shaking(best, k+1, n, W, w)
            x_second = neighborhoods[k](x_prime, n, W, u, w)
            
            if get_weight(x_second, w) <= W and get_value(x_second, u) > get_value(best, u):
                best = x_second
                k = 0 
            else:
                k += 1 
    
    return best, time.perf_counter() - start
\end{lstlisting}

\end{document}